%----------------------------------------------------------------------------------------
%	COVER LETTER
%----------------------------------------------------------------------------------------

% To remove the cover letter, comment out this entire block
\clearpage

 \newenvironment{indented}{%
    \par % Start a new paragraph
    \setlength{\parindent}{2em} % Set your preferred indent size
    \indent % Force the indent on the very first line
}{%
    \par % Ensure the last paragraph finishes before closing
}
\recipient{}{} % Letter recipient
\date{\today} % Letter date
\opening{\textbf{Summer Funding Reserch Plan}} % Opening greeting
\closing{} % Closing phrase
\enclosure[Attached]{Curriculum Vit\ae{}} % List of enclosed documents

% --- START MINI-TITLE PAGE ---
\begin{center}
    % Removes page number from the title page
    
    \vspace*{4cm} % Pushes the content down the page
    
    {\Huge \textbf{Hamilton Award Application}} \\ [1.5cm]
    {\Large Alexander Marshall} \\ [0.5cm]
    {\large \today}
    
    \vfill % Pushes everything else to the top, leaving the bottom empty
\end{center}
% --- END MINI-TITLE PAGE ---

% --- START MANUAL TOC ---
\begin{center}
    {\Large \textbf{Table of Contents}}
\end{center}

\vspace{1cm}

\noindent \textbf{1. Other Sources of Funding} \dotfill \textbf{1}\\
\noindent \textbf{2. Summer Funding Plan} \dotfill \textbf{2} \\
\noindent \textbf{3. Letter of Recommendation} \dotfill \textbf{3}\\
\noindent \textbf{4. Curriculum Vit\ae} \dotfill \textbf{4}\\ 
\par \vspace{2cm}
\textbf{Other Sources of Funding}\par
I am a Royster Doctoral Fellow. Presently, I am scheduled to spend the Summer of 2026 as a non-service semester. If I do not receive the Hamilton award, I will receive \$5000 as a stipend from the Graduate School. The Royster Fellowship is flexible, and if I do receive the Hamilton Award, I will some amount of this funding to a later date so I can receive the full benefits of the Award.

\newpage
\textbf{Summer Funding Plan}\par

\begin{indented}
I intend to use the Hamilton Award to enable full-time commitment to research this summer. Doing so will allow me to focus the entirety of my energies to research, speeding progress by enabling me to advance my research goals more quickly. \par 

Tissues and individual cells are poroelastic, consisting of a background matrix suffused by a solvent carrying solutes of varying sizes. In vivo, materials undergo dynamic stresses, causing deformations which induce advection and size-dependent transport. Solute distribution in the nucleus, a heterogeneous poroelastic environment often subjected to dynamic strains, drives gene transcription. Experiments examining the macroscale dynamics of poroelastic media demonstrate enhanced concentration of large solutes in cyclically strained hydrogels, but there remains a dearth of studies probing the intricate interactions between solutes, pores, stress gradients, and microfluidics during dynamic strains, and thus the roots of these emergent behaviors remain unclear. Diffusive transport of large solutes in poroelastic media is anomalous, strain-dependent, and temporally anisotropic, necessitating high spatiotemporal resolution. \par

My work to date has consisted of tracking single nanoparticles in strained poroelastic media, providing the first live insight into the nanoscopic dynamics of these systems. I have shown that both static and dynamic strains increase long-term localization and decrease transport compared to unstrained controls, dynamic strain induces asymmetric advective transport of large solutes during compression and relaxation, creating a pumping effect. All of the work to date has been completed with the eventual goal of conducting similar work in live-cell imaging.\par

This summer, I will extend the study to living systems by tracking transcription factors (TFs) in a strained nuclear envelope to understand how strains impact the transport of transcription factors within living organisms. Doing so will require fluorescently labelled transcription factors to be imaged in live cellular nuclei subject to controlled strains.\par
To accomplish this, human macrophages will be cultured and treated with a non-toxic HALO-tag solution to fluorescently labelled transcription factors within the nuclear envelope. The cells will be imaged on a light sheet fluorescence microscopy system which operates with low phototoxicity inside a climate-controlled incubation system, allowing for extended experimental sessions on healthy cells. Precise strain will be applied during live imaging using a 20$\mu m$ spherical indenter mounted on a piezoelectric-actuated atomic force microscope cantilever. The strain system is mounted on a micro-manipulation stage, allowing for precise indenter positioning.
High speed imaging of fluorescently-labelled TFs will be conducted using light sheet fluorescence microscopy to acquire high sensitivity, high-framerate time series of the particles. Particle tracking will be conducted using TrackMate, open-source state-of-the-art particle tracking software, and the data will be analyzed using MATLAB scripts. As an extension, an instrument unique to the Superfine Lab which enables simultaneous atomic force microscopy and light sheet fluorescence microscopy may be used to produce simultaneous particle tracking and microrheological data.\par 

The overall approach detailed above been used in my prior gel-based strain experiments to successfully track comparable nanoparticles in strained poroelastic media at up to 800fps. Furthermore, recent collaborations with Dr. Kerry Bloom of UNC Biology have demonstrated single particle tracking of HALO-tagged histone proteins in the nucleus. These prior experiments were designed in part as milestones along the pathway towards this experiment detailed in this proposal. Years have been spent developing a system with the capabilities necessary for this work, and recent experiments demonstrate the achievement of anticipated functionality requirements.
\end{indented}
\newpage
\textbf{Letter of Recommendation}\par

\begin{indented}
    Alex is so Great!
\end{indented}

\newpage
